\section{Related Work}
\label{sec:related}
% Rewritten 2026-09-17 into the user's three subsections. Citations are limited to entries that
% exist in the bibliography and were verified; where a canonical reference is missing it is named in a
% comment rather than invented. `ase26' is pre-existing and still has no entry.

\subsection{Comprehension Under Obfuscation and Deobfuscation}
Obfuscation has a standard taxonomy of lexical, control and data
transforms~\cite{collberg1997taxonomy,collberg2002manufacturing}, and its effect on \emph{human}
comprehension is measured: Nguyen et al.\ show it significantly reduces analysts' accuracy at
predicting program outputs~\cite{nguyen2026effectcodeobfuscationhuman}. For language models, recent
work reports the same degradation under semantics-preserving transforms, including on output
prediction~\cite{ase26}, and a decline that grows with obfuscation
complexity~\cite{nikiema2025codebarrierllmsactually}. Our evaluation extends that line in two ways:
composition, since prior work evaluates single transforms and real obfuscators stack them, and
direction, since we ask the reverse question on the same programs.

We deliberately do not join the deobfuscation lineage, in which recovering the original source from
its obfuscated form is the training objective itself~\cite{dobf}. That is a different question: a model that can undo a transform
tells us about the transform, whereas a model that answers correctly on the transformed program tells
us about its understanding of the computation. Every evaluation in this paper is on code that stays
obfuscated.

\subsection{Model Merging and Modular Adaptation}
Low-rank adaptation~\cite{lora} made per-task adapters cheap enough that combining them became a
design question in its own right. Weight-space merging answers it without further training. Task arithmetic showed that the difference
between fine-tuned and pre-trained weights is a vector that can be added and negated to steer
behaviour~\cite{taskarithmetic}; TIES
resolves interference between task vectors by trimming and sign election~\cite{ties}, and DARE shows
that fine-tuning deltas are redundant enough to drop most entries and rescale the
rest~\cite{dare}. Both are used here unmodified. Our contribution to this literature is not a
merging method but a comparison of \emph{where} merging sits relative to pooled training and routing
under distribution shift --- and the finding that its advantage appears precisely where the training
distribution stops, which the merging literature, evaluated in-distribution, does not exercise.

\subsection{Direction Specificity in Code Models}
Benchmarks that pose both input and output prediction on the same functions established that current
code models are imperfect at both and that the two are not interchangeable~\cite{cruxeval}. We use
that pairing as an instrument rather than a benchmark: because the two directions share a program and
a gold value, an adaptation that improves one and damages the other has changed what it produces
rather than what it understands. The forward-collapse measurement --- the rate at which a model asked
for an input returns the output --- makes that specificity directly observable, and to our knowledge
it has not previously been used to separate task fitting from program understanding in adapted code
models.
